\documentclass[11pt,a4paper]{article}
\usepackage{kotex}
\usepackage{fontspec}
\setmainfont{Times New Roman}
\setmainhangulfont{AppleMyungjo}
\setsanshangulfont{Apple SD Gothic Neo}
\usepackage[a4paper,top=18mm,bottom=18mm,left=20mm,right=20mm]{geometry}
\usepackage{fancyhdr}
\setlength{\headheight}{24pt}
\usepackage{enumitem}
\usepackage{mdframed}
\usepackage{booktabs}
\linespread{1.18}
\setlength{\parindent}{0pt}
\pagestyle{fancy}
\fancyhf{}
\renewcommand{\headrulewidth}{0.6pt}
\fancyhead[L]{\sffamily\bfseries 2026 고1 영어 변형 - 정답 및 해설 지문 33}
\fancyhead[R]{\sffamily 후각 묘사의 어려움}
\fancyfoot[C]{\framebox[16mm][c]{\thepage}}
\newcommand{\qno}[1]{\par\vspace{10pt}{\sffamily\bfseries #1.}\ }
\newmdenv[linewidth=0.6pt,roundcorner=3pt,innertopmargin=7pt,innerbottommargin=7pt,
          innerleftmargin=9pt,innerrightmargin=9pt,skipabove=5pt,skipbelow=5pt]{pbox}
\begin{document}
\begin{center}
{\fontsize{15}{17}\selectfont\sffamily\bfseries [지문 33] 정답 및 해설}
\end{center}
\vspace{6pt}
\begin{center}
\renewcommand{\arraystretch}{1.35}
\begin{tabular}{cccc}
\toprule
문항 & 유형 & 정답 & 배점 \\
\midrule
1 & 빈칸추론 & ③ & 3점 \\
2 & 어휘 부적절 & ③ & 2점 \\
3 & 서술형 단답 & 모범답안 참조 & 2점 \\
4 & 요약문 & ① & 2점 \\
\bottomrule
\end{tabular}
\end{center}

\qno{1}\textbf{빈칸추론 - 정답: ③}
\begin{pbox}
빈칸 뒤의 예시 ``flowery,'' ``fruity,'' ``fishy''는 냄새 자체를 독립적인 형용사로 묘사하는 것이 아니라, 그 냄새를 만들어 내는 원천을 빌려 말하는 방식이다. 따라서 ``the source of the odor''와 같은 뜻인 ③ \textit{what produces the smell}이 정답이다.

\medskip
① pleasantness, ② identification ability, ④ emotional reaction, ⑤ familiarity는 모두 냄새에 대해 말할 수 있는 요소처럼 보이지만, 바로 뒤 예시들이 꽃, 과일, 생선이라는 ``냄새의 원천''을 가리킨다는 점과 맞지 않는다.
\end{pbox}

\qno{2}\textbf{어휘 부적절 - 정답: ③}
\begin{pbox}
본문의 핵심은 ``매우 익숙한 냄새조차도 일상적 맥락 밖에서는 이름 붙이기 어렵다''는 것이다. 따라서 ③ \textit{unfamiliar}는 문맥상 부적절하며 \textit{familiar}가 되어야 한다.

\medskip
① challenging은 후각 경험을 말로 설명하기 어렵다는 내용과 맞고, ② identify는 냄새를 알아보고 이름 붙이는 실험 내용과 맞다. ④ unable은 참가자 대부분이 흔한 냄새의 이름을 대지 못했다는 결과와 맞고, ⑤ shaped는 진화가 우리의 뇌를 그렇게 만들었다는 결론과 맞다.
\end{pbox}

\qno{3}\textbf{서술형 단답}
\begin{pbox}
\textbf{모범답안}: 대부분의 참가자가 흔한 냄새의 이름을 대지 못했다.

\medskip
\textbf{허용답안}: 일상적인 냄새조차 맥락 밖에서는 잘 식별하지 못했다. / 참가자 다수가 맥주, 장미 같은 흔한 냄새를 이름 붙이지 못했다.

\medskip
\textbf{채점 기준}: ``대부분의 참가자'' 또는 이에 준하는 다수성, ``흔한/익숙한 냄새'', ``이름을 대지 못함/식별하지 못함''이 포함되면 정답.
\end{pbox}
\vspace{8pt}
\qno{4}\textbf{요약문 - 정답: ①}
\begin{pbox}
글의 결론은 인간이 냄새에 실제 상황에서는 올바르게 반응할 수 있지만, 그것을 말로 설명하거나 이름 붙이는 데에는 서툴다는 것이다. 따라서 (A)는 \textit{correctly}, (B)는 \textit{describing}이 적절하다.
\end{pbox}

\end{document}
